% GENERATED FILE. Edit canonical lessons, then run npm run generate:books.
% canonical-source-hash: fnv1a64:eb35bf8f26b42acd
% canonical-lessons: MR-C136-kot, MR-C136-svetar, MR-C136-poli, MR-C136-bhaji, MR-C136-varan

\chapter{Coat, Jumper, Chapati, Vegetable dish, Plain dal}
\label{ch:mr-coat-jumper-chapati-vegetable-dish-plain-dal}
\hlchaptermodality{136}

\begin{chapteropening}
\textbf{By the end of this chapter:} \emph{I can say a coat, a jumper, a chapati, a vegetable dish and plain dal.}

Complete the last lesson of chapter 136: I can say a coat, a jumper, a chapati, a vegetable dish and plain dal.
\end{chapteropening}

\section[\texorpdfstring{\mr{कोट} (koṭ)}{कोट (koṭ)}]{\mr{कोट} (koṭ) — a coat}

\label{lesson:MR-C136-kot}



\begin{quote}
Before the new one: say the Marathi for study, then the Marathi for a shirt.
\end{quote}

\subsection*{You'll want to know: \mr{कोट}}
\textbf{\mr{कोट}} — \emph{koṭ} — \textquotedblleft{}a coat\textquotedblright{}.

\begin{cousinweb}
\textbf{\mr{कोट}} (\emph{koṭ}) is a loan from English \emph{coat}. It is not the inherited Marathi \emph{koṭ}, \textquotedblleft{}fort\textquotedblright{}.
\end{cousinweb}

\begin{grammarlens}[title={What you've built}]
One more word for this chapter.
\end{grammarlens}

\subsection*{Your turn}
\begin{itemize}
\raggedright
  \item \emph{Say it:} \emph{koṭ}
  \item \emph{Say it:} \emph{koṭ}, once more
  \item \emph{Say it:} say \emph{koṭ}
  \item \emph{Recall:} say the Marathi for study, then the Marathi for a shirt, then say \emph{koṭ} again
\end{itemize}

\begin{tcolorbox}[breakable,colback=teal!4,colframe=teal!35!black,title={Before you move on}]
What does \mr{कोट} mean? (\textquotedblleft{}A coat\textquotedblright{}.) Say it once more.
\end{tcolorbox}
\section[\texorpdfstring{\mr{स्वेटर} (sveṭar)}{स्वेटर (sveṭar)}]{\mr{स्वेटर} (sveṭar) — a jumper}

\label{lesson:MR-C136-svetar}



\begin{quote}
Before the new one: say the Marathi for a shirt, then the Marathi for a coat.
\end{quote}

\subsection*{You'll want to know: \mr{स्वेटर}}
\textbf{\mr{स्वेटर}} — \emph{sveṭar} — \textquotedblleft{}a jumper\textquotedblright{}.

\begin{cousinweb}
\textbf{\mr{स्वेटर}} (\emph{sveṭar}) is a loan from English \emph{sweater}.
\end{cousinweb}

\begin{grammarlens}[title={What you've built}]
One more word for this chapter.
\end{grammarlens}

\subsection*{Your turn}
\begin{itemize}
\raggedright
  \item \emph{Say it:} \emph{sveṭar}
  \item \emph{Say it:} \emph{sveṭar}, once more
  \item \emph{Say it:} say \emph{sveṭar}
  \item \emph{Recall:} say the Marathi for a shirt, then the Marathi for a coat, then say \emph{sveṭar} again
\end{itemize}

\begin{tcolorbox}[breakable,colback=teal!4,colframe=teal!35!black,title={Before you move on}]
What does \mr{स्वेटर} mean? (\textquotedblleft{}A jumper\textquotedblright{}.) Say it once more.
\end{tcolorbox}
\section[\texorpdfstring{\mr{पोळी} (poḷī)}{पोळी (poḷī)}]{\mr{पोळी} (poḷī) — a chapati}

\label{lesson:MR-C136-poli}



\begin{quote}
Before the new one: say the Marathi for a coat, then the Marathi for a jumper.
\end{quote}

\subsection*{You'll want to know: \mr{पोळी}}
\textbf{\mr{पोळी}} — \emph{poḷī} — \textquotedblleft{}a chapati\textquotedblright{}.

\begin{cousinweb}
\textbf{\mr{पोळी}} (\emph{poḷī}) is a tadbhava, a word inherited from Sanskrit and reshaped by sound change. It comes from \emph{pauli}, \textquotedblleft{}a cake of scorched grain and ghee\textquotedblright{} (CDIAL 8423). CDIAL glosses the Marathi word \textquotedblleft{}wheaten cake\textquotedblright{}, and suggests, with a \textquotedblleft{}perhaps\textquotedblright{}, a link with Prakrit \emph{paülaï}, \textquotedblleft{}cooks\textquotedblright{}, or with \emph{pūpalā}, \textquotedblleft{}a cake\textquotedblright{}.
\end{cousinweb}

\begin{grammarlens}[title={What you've built}]
One more word for this chapter.
\end{grammarlens}

\subsection*{Your turn}
\begin{itemize}
\raggedright
  \item \emph{Say it:} \emph{poḷī}
  \item \emph{Say it:} \emph{poḷī}, once more
  \item \emph{Say it:} say \emph{poḷī}
  \item \emph{Recall:} say the Marathi for a coat, then the Marathi for a jumper, then say \emph{poḷī} again
\end{itemize}

\begin{tcolorbox}[breakable,colback=teal!4,colframe=teal!35!black,title={Before you move on}]
What does \mr{पोळी} mean? (\textquotedblleft{}A chapati\textquotedblright{}.) Say it once more.
\end{tcolorbox}
\section[\texorpdfstring{\mr{भाजी} (bhājī)}{भाजी (bhājī)}]{\mr{भाजी} (bhājī) — a vegetable dish}

\label{lesson:MR-C136-bhaji}



\begin{quote}
Before the new one: say the Marathi for a jumper, then the Marathi for a chapati.
\end{quote}

\subsection*{You'll want to know: \mr{भाजी}}
\textbf{\mr{भाजी}} — \emph{bhājī} — \textquotedblleft{}a vegetable dish\textquotedblright{}.

\begin{cousinweb}
\textbf{\mr{भाजी}} (\emph{bhājī}) is a tadbhava, a word inherited from Sanskrit and reshaped by sound change. It comes from \emph{bharjita}, \textquotedblleft{}fried, parched\textquotedblright{} (CDIAL 9401).
\end{cousinweb}

\begin{grammarlens}[title={What you've built}]
One more word for this chapter.
\end{grammarlens}

\subsection*{Your turn}
\begin{itemize}
\raggedright
  \item \emph{Say it:} \emph{bhājī}
  \item \emph{Say it:} \emph{bhājī}, once more
  \item \emph{Say it:} say \emph{bhājī}
  \item \emph{Recall:} say the Marathi for a jumper, then the Marathi for a chapati, then say \emph{bhājī} again
\end{itemize}

\begin{tcolorbox}[breakable,colback=teal!4,colframe=teal!35!black,title={Before you move on}]
What does \mr{भाजी} mean? (\textquotedblleft{}A vegetable dish\textquotedblright{}.) Say it once more.
\end{tcolorbox}
\section[\texorpdfstring{\mr{वरण} (varaṇ)}{वरण (varaṇ)}]{\mr{वरण} (varaṇ) — plain dal}

\label{lesson:MR-C136-varan}



\begin{quote}
Before the new one: say the Marathi for a chapati, then the Marathi for a vegetable dish.
\end{quote}

\subsection*{You'll want to know: \mr{वरण}}
\textbf{\mr{वरण}} — \emph{varaṇ} — \textquotedblleft{}plain dal\textquotedblright{}.

\begin{cousinweb}
CDIAL lists Marathi \emph{varaṇ}, \textquotedblleft{}tasty dish of pulse\textquotedblright{}, under a reconstructed \emph{*aparānna}, built from Sanskrit \emph{apara}, \textquotedblleft{}other, later\textquotedblright{}, and \emph{anna}, \textquotedblleft{}food\textquotedblright{} (CDIAL 442). CDIAL calls the link very doubtful: the loss of the first \emph{a-} and the single \emph{-ṇ-} are not explained.
\end{cousinweb}

\begin{grammarlens}[title={What you've built}]
One more word for this chapter.
\end{grammarlens}

\subsection*{Your turn}
\begin{itemize}
\raggedright
  \item \emph{Say it:} \emph{varaṇ}
  \item \emph{Say it:} \emph{varaṇ}, once more
  \item \emph{Say it:} say \emph{varaṇ}
  \item \emph{Recall:} say the Marathi for a chapati, then the Marathi for a vegetable dish, then say \emph{varaṇ} again
\end{itemize}

\begin{tcolorbox}[breakable,colback=teal!4,colframe=teal!35!black,title={Before you move on}]
What does \mr{वरण} mean? (\textquotedblleft{}Plain dal\textquotedblright{}.) Say it once more.
\end{tcolorbox}
